\documentclass[book.tex]{subfiles}
\begin{document}

\chapter{Molecular Graph Convolutions}

\label{chap:molecular_graphconvs}
\noindent\rule{11cm}{0.4pt} \\
\textbf{Prerequisites:} \autoref{chap:graphconvs} \\
\textbf{Difficulty Level:} ** \\
\noindent\rule{11cm}{0.4pt}
\vspace{5mm}

Molecular graph convolutions treat a molecule as a graph and utilize graph convolutional techniques to make predictions about molecular properties. These convolutions provide a powerful alternative to the heuristic fingerprints we learned about earlier since they allow for the possibility of learning data-dependent featurizations suited to the problem at hand instead of using generic featurization schemes. Empirically, molecular graph convolutions can outperform heuristic featurizations on a broad range of tasks.

\section{Representing Molecules as Graphs}
Graph-based models are naturally suited to represent molecules as mathematical graphs by defining atoms as nodes, bonds as edges. The natural similarity has inspired a number of models to utilize the graph structure of molecules to improve predictive capabilities\cite{gilmer2017neural}. Mathematically, the representation of a molecule is given by vertices $\mathcal{V}$ and edges $\mathcal{E}$.

%\textcolor{red}{TODO: Add an introduction to graphs in the mathematical section.}
\begin{align}
\mathcal{V} &= \{1, \dotsc, n\}\\
\mathcal{E} &= \{(i, j)\ | \textrm{atom $i$ is covalently bonded to atom $j$}\} \\
\end{align}

Here $n$ is the number of atoms in the molecule, $\mathcal{V}$ represents the set of vertices (atoms) and $\mathcal{E}$ the set of edges (covalent bonds). The molecule is represented by the pair
\begin{align}
m &= (\mathcal{V}, \mathcal{E})
\end{align}
We introduce a graph data type in Physika to facilitate coding
\begin{lstlisting}
class Node(h: $\mathbb{R}^N$)
class Graph(V: [Node], E: [(Node, Node)])
\end{lstlisting}
As a first summary, graph convolutional models treat molecules as undirected graphs, and apply the same learnable function to every node (atom) and its neighbors (bonded atoms) in the graph.

\section{Neural Graph Fingerprints}

A neural graph fingerprint architecture begins with computing an initial feature vector and a neighbor list for every atom in the molecule. The feature vector encodes the atom's local electronic structure or the chemical environment, by incorporating atom-types, hybridization types, and valence structures. The neighbor lists of atoms represent connectivity within the molecule, which are further processed in each model to generate graph structures.

\begin{align}
    \textrm{Features} &\in \mathbb{R}^{N \times d} \\
    \textrm{NeighborList} &= [[i_1,\dotsc,i_j],\dotsc]
\end{align}

As for ECFP fingeprints, graph convolutional layer gradually merge information for distant atoms by extending radially through bonds. However, instead of applying fixed hash functions, the neural graph fingerprint uses differentiable "hash functions". This creates a learnable process capable of extracting useful representations of molecules suited to the task at hand, which recapitulates convolution layers in visual recognition deep networks.

\begin{figure}
    \centering
    % \includegraphics{figures/Differentiable Physics/Molecular ML/Molecular Graph Convolutions/neural_graph_fingerprint.png}
    \includegraphics{figures/Differentiable Physics/Molecular ML/Molecular Graph Convolutions/neural_graph_fingerprint0.png}
    \caption{An illustration of the control flow in a classical molecular fingerprint and a neural graph fingerprint. \url{https://arxiv.org/pdf/1509.09292.pdf}}
    \label{fig:neural_graph_fingerprint}
\end{figure}

We can implement the neural graph fingerprint in Physika as follows. Let \lstinline{R} denote the number of layers in the architecture. Let $H^i_j$ denote the weight matrix for the $j$-th hidden layer and the case when an atom has $i$ neighbors, where $N$ is the max number neighbors we allow (typically 5).

\begin{lstlisting}[language=python]
class NeuralGraphFingerprint(R: $\mathbb{N}$, [$H^1_1$,$\dotsc$,$H^N_R$], $W_1,\dotsc, W_R$):

  def $\lambda$(molecule: Graph):
    for a in molecule:
      r[a] = g(a)
    for L in [1,$\dotsc$, R]:
      for a in molecule:
        $r_1$,$\dotsc$,$r_N$ = neighbors(a)
        v = r[a] + $\sum_{i=1}^N r_i$
        $r_a$ = $\sigma$(v$H^N_L$)
        i = softmax($r_a$ $W_L$)
        f = f + i
    return f
\end{lstlisting}

Note the similarity between the neural graph fingerprint and the ECFP algorithm. The core difference is the presence of the learnable matrices $H^i_j$ which allow the emphasis the fingerprint places on different parts of the molecule to be tuned by data.

\section{Weave Convolutions}

The weave featurization \cite{kearnes2016molecular} is a modification of the graph convolutional featurization that captures more bond information. The atom feature vector in the weave featurization is exactly the same as in the graph convolutional featurization, while the connectivity between atoms is represented using more detailed pair features instead of a neighbor listing. 

For each pair of atoms in the molecule, the weave featurization calculates a feature vector, including bond properties if atoms are directly bonded, distance and ring information, creating a feature matrix. This featurization method can be used in conjunction with graph-based models that utilize properties of both nodes (atoms) and edges (bonds between atoms).

Unlike the neural graph fingerprint, the weave convolution is much larger; to update features of an atom, weave models combine info from all other atoms and their corresponding pairs in the molecule. Weave models are more efficient at transmitting information between distant atoms, at the price of increased complexity for each convolution. A molecule is first encoded into a list of atomic features and a matrix of pair features by the weave model's featurization method. Then in each weave module, these features are inputted into four sets of fully connected layers (corresponding to four paths from two original features to two updated features) and concatenated to form new atomic and pair features. After stacking several weave modules, a similar gather layer combines atomic features together to form molecular features that are fed into task-specific layers.

Let $x_1,\dotsc x_n$ be layers. Atom layers can be computed from previous atom layers ($A \mapsto A$) by applying

\begin{align}
    A^y_a &= f(A^{x_1}_a,\dotsc,A^{x_n}_a)
\end{align}

Pair layers can be computed from previous pair layers ($P\mapsto P$) by applying

\begin{align}
    P^y_{a,b} &= f(P^{x_1}_{a,b}, \dotsc, P^{x_n}_{a,b})
\end{align}

We can also update atom layers from pair layers ($P \mapsto A$) by considering all pairs which contain a given atom $a$.
\begin{align}
    A^y_a &= g(f(P^x_{(a,b)}), f(P^x_{(a,c)}),\dotsc )
\end{align}

\begin{figure}
    \centering
    % \includegraphics[width=0.5\textwidth]{figures/Differentiable Physics/Molecular ML/Molecular Graph Convolutions/P_to_A_weave.png}
    \includegraphics[width=0.6\textwidth]{figures/Differentiable Physics/Molecular ML/Molecular Graph Convolutions/Pair to atom.png}
    \caption{The $P \mapsto A$ transformation projects from pair features to atom features. \url{https://arxiv.org/abs/1603.00856}}
    \label{fig:p_to_a}
\end{figure}

We can also expand a pair layer from an atom layer ($A \mapsto P$).
\begin{align}
    P^y_{a,b} &= g(f(A^x_a, A^x_b), f(A^x_b, A^x_a))
\end{align}

\begin{figure}
    \centering
    % \includegraphics[width=0.5\textwidth]{figures/Differentiable Physics/Molecular ML/Molecular Graph Convolutions/A_to_P_weave.png}
    \includegraphics[width=0.5\textwidth]{figures/Differentiable Physics/Molecular ML/Molecular Graph Convolutions/Atom to pair.png}
    \caption{The $A \mapsto P$ transformation combines atom features to generate pair features. \url{https://arxiv.org/abs/1603.00856}}
    \label{fig:my_label}
\end{figure}

We can combine all these operations together into a Weave module

\begin{figure}
    \centering
    % \includegraphics[width=0.5\textwidth]{figures/Differentiable Physics/Molecular ML/Molecular Graph Convolutions/weave_module.png}
    \includegraphics[width=0.5\textwidth]{figures/Differentiable Physics/Molecular ML/Molecular Graph Convolutions/weave module.png}
    \caption{The Weave module combines atom and pair features to generate new atom and pair features by performing $A \mapsto A$, $P \mapsto P$, $A \mapsto P$, and $P \mapsto A$ transformations. \url{https://arxiv.org/abs/1603.00856}}
    \label{fig:my_label}
\end{figure}

We can represent the Weave module in Physika as follows.

\begin{lstlisting}[language=python]
class Weave(A $\mapsto$ A: $\mathbb{R}[N, d_a] \to \mathbb{R}[N, d_a]$, A $\mapsto$ P: $\mathbb{R}[N, d_a] \to \mathbb{R}[N, N, d_p]$, 
            P $\mapsto$ A: $\mathbb{R}[N, N, d_p] \to \mathbb{R}[N, d_a]$, P $\mapsto$ P: $\mathbb{R}[N, N, d_p] \to \mathbb{R}[N, N, d_p]$):
  def $\lambda$($A^k$: $\mathbb{R}[N, d_a]$, $P^k$: $\mathbb{R}[N, N, d_p]$):
    $A^{k'}$ = ($A\mapsto A$)($A^k$)
    $A^{k''}$ = ($P \mapsto A$)($P^k$)
    $A^{k+1}$ = ($A \mapsto A$)($A^{k'}$ || $A^{k''}$)
  
    $P^{k'}$ = ($P\mapsto P$)($P^k$)
    $P^{k''}$ = ($A \mapsto P$)($A^k$)
    $P^{k+1}$ = ($P \mapsto P$)($P^{k'}$ || $P^{k''}$)
    return $A^{k+1}, P^{k+1}$
\end{lstlisting}

Weave convolutions can be stacked. At the end of a stack of weave convolutions, the final atom features $A^N$ can be extracted and used in a downstream pipeline.

\begin{figure}
    \centering
    % \includegraphics[width=0.5\textwidth]{figures/Differentiable Physics/Molecular ML/Molecular Graph Convolutions/weave_architecture.png}
    \includegraphics[width=0.3\textwidth]{figures/Differentiable Physics/Molecular ML/Molecular Graph Convolutions/Stacked weave module.png}
    \caption{A stack of weave modules can be combined to build a feature vector that can be fed into a fully connected network. \url{https://arxiv.org/abs/1603.00856}}
    \label{fig:weave_arcvh}
\end{figure}

% TODO: Maybe add this in? Just not sure it's actually useful
%\subsection{Directed Acyclic Graph Models}
%(refine text) Directed Acyclic Graph (DAG) models regard molecules as directed graphs. While chemical bonds typically do not have natural directions, one can arbitrarily generate a DAG on a molecule by designating a central atom and then define directions of all bonds in certain orientations towards the atom.\cite{lusci2013deep} In the case of small molecules, taking all possible orientations is computationally feasible. In other words, for a molecule with $n_a$ atoms, the model will generate $n_a$ DAGs, each centered on a different atom. In the actual calculations of a graph, a vector of graph features is calculated for each atom based on its atomic features (reusing the graph convolutions featurizer) and its parents' graph features. As features gradually propagate through bonds, information converges on the central atom. Then a final sum of all graphs gives the molecular features, which are fed into classification or regression tasks. Note that $n_a$ graphs are evaluated for each molecule, which can cause a significant increase in required calculations.


\end{document}